\documentclass[a4paper,12pt,twoside,openright]{mgr}

%**************************************************************************
% Dane do strony tytułowej
%

% Autor
\autor{inż. Dawid Ziora}

% Rodzaj pracy - wpisać LICENCJACKA, INŻYNIERSKA lub MAGISTERSKA
\rodzajPracy{MAGISTERSKA}

% Tytuł pracy magisterskiej/inżynierskiej
\tytul{Wpływ wybranych frameworków Java na zwięzłość i czytelność kodu źródłowego}

% Tytuł pracy magisterskiej w języku angielskim
\tytulAng{The impact of selected Java frameworks on the conciseness and readability of source code}

% Promotor
\promotor{dr inż. Grzegorz Michalski}

% Rok
\rok{2026}

% Kierunek
\kierunek{Informatyka}

% Studia stacjonarne lub niestacjonarne (wpisać jakie)
\studia{stacjonarne}

% Poziom studiów wpisać I lub II
\poziomStudiow{II}

% Numer albumu
\numerAlbumu{136700}

%
%**************************************************************************

% Styl dla wtrąceń anglojęzycznych
\newcommand{\eng}[1]{(\emph{#1})}
%\renewcommand{\baselinestretch}{1}
\renewcommand\lstlistlistingname{Spis listingów}

\begin{document}

\pagestyle{empty}
\stronaTytulowa
\cleardoublepage
\tableofcontents
\pagestyle{fancy}

\chapter{Wstęp}
W pracy ... \cite{AcademyBank}
\section{Cel pracy}
\section{Zakres badań}
\section{Struktura pracy}
\chapter{Przegląd technologii}
\section{Java}
\subsection*{Quarkus}
\subsection*{Spring Boot}
\chapter{Narzędzia}
\chapter{Analiza wyników badań}
\chapter{Podsumowanie i wnioski}
\chapter*{Streszczenie}
\addcontentsline{toc}{chapter}{Streszczenie}
\chapter*{Abstract}
\addcontentsline{toc}{chapter}{Abstract}
\printbibliography[heading=bibintoc]

\listoffigures

\listoftables

\lstlistoflistings

\addcontentsline{toc}{chapter}{\lstlistlistingname}
\end{document}
